\section{Funktionale Abhängigkeiten (FA)}
\subsection{Funktionsbegriff, reelle Funktionen, Darstellungen und Eigenschaften}
Eine Funktion ist eine Zuordnung, die jedem Element $x$ aus der Definitionsmenge $D$ eine Element $y$ aus der Wertemenge $W$ zuordnet. Eine Funktion wird üblicherweise folgendermaßen angeschrieben:
\begin{equation}
    f: D \rightarrow W,\; y=f(x)
\end{equation}
Der Begriff Abbildung ist ebenfalls ein gebräuchlicher Name für Funktionen. Oft werden Funktionen mit dem Abbildungspfeil $\mapsto$ dargestellt:
\begin{equation}
    f: D \rightarrow W,\; x \mapsto f(x)
\end{equation}
Wichtig ist auch zu unterscheiden, ob es sich tatsächlich um eine Funktion handelt. Hierbei ist wichtig zu beachten, dass es zu jedem Element $x$ genau ein Element $y$ gibt. Zur Prüfung kann man eine senkrechte Linie verwenden.
\vspace{10pt}
\\
\noindent
\begin{minipage}[t]{.33\textwidth}
\begin{center}
\begin{tikzpicture}[scale=1.2, >=stealth]
    \draw[->, thick] (0,0) -- (3,0) node[right] {$x$};
    \draw[->, thick] (0,0) -- (0,2) node[above] {$y$};
    % Kurvige Funktion (z.B. kubische Funktion)
    \draw[accent, ultra thick, domain=0:3, smooth, samples=100]
        plot (\x, {0.2*(\x-1.5)^3 + 1});
\end{tikzpicture}
\captionof*{figure}{eine Funktion}
\end{center}
\end{minipage}
\begin{minipage}[t]{.33\textwidth}
\begin{center}
\begin{tikzpicture}[scale=1.2, >=stealth]
    \draw[->, thick] (0,0) -- (3,0) node[right] {$x$};
    \draw[->, thick] (0,0) -- (0,2) node[above] {$y$};
    \draw[accent, ultra thick, domain=0:360, smooth, samples=200]
        plot ({1.5 + 0.9*cos(\x)}, {1 + 0.7*sin(\x)});
    \draw[dashed, thick] (1.5,0) -- (1.5,1.9) node[above]{};
    \fill (1.5,0.3) circle (1.5pt) node[above=2pt] {};
    \fill (1.5,1.7) circle (1.5pt) node[above=2pt] {};
\end{tikzpicture}
\captionof*{figure}{keine Funktion}
\end{center}
\end{minipage}
\begin{minipage}[t]{.33\textwidth}
\begin{center}
\begin{tikzpicture}[scale=1.2, >=stealth]
    \draw[->, thick] (0,0) -- (3,0) node[right] {$x$};
    \draw[->, thick] (0,0) -- (0,2) node[above] {$y$};
    \draw[accent, ultra thick, domain=0:360, smooth, samples=200]
        plot ({1.5 + 0.9*cos(\x)}, {1 + 0.4*sin(2*\x)});
    \draw[dashed, thick] (1,0) -- (1,1.9) node[above]{};
    \fill (1,0.63) circle (1.5pt) node[above=2pt] {};
    \fill (1,1.37) circle (1.5pt) node[above=2pt] {};
\end{tikzpicture}
\captionof*{figure}{keine Funktion}
\end{center}
\end{minipage}
\vspace{10pt} \\
\noindent
Möchte man den Graph einer Funktion skizzieren, kann es hilfreich sein eine Wertetabelle anzufertigen.
\begin{bsp}
    Darstellungsarten von reellen Funktionen: \\
    Wir untersuchen folgende Funktion:
    \[
        f(x)=0{,}5x^2+1
    \]
    Zuerst wird eine Wertetabelle angefertigt und daraufhin der Graph der Funktion dargestellt.
    \begin{center}
        \begin{minipage}{0.4\textwidth}
            \centering
            \begin{tabular}{c|c}
                $x$ & $f(x)$ \\ \hline
                $-2$ & $3$ \\
                $-1$ & $1.5$ \\
                $0$  & $1$ \\
                $1$  & $1.5$ \\
                $2$  & $3$
            \end{tabular}
        \end{minipage}
        \hfill
        \begin{minipage}{0.55\textwidth}
            \centering
            \begin{tikzpicture}
                \begin{axis}[
                    axis lines = middle,
                    xmin = -3, xmax = 3,
                    ymin = 0,  ymax = 5,
                    xtick={-2,-1,0,1,2},
                    ytick={0,1,2,3,4},
                    samples = 200,
                    % grid = major,
                    width = 7cm,
                    height = 5cm,
                    xlabel = {$x$},
                    ylabel = {$y$},
                    domain = -3:3
                ]
                \addplot[accent, ultra thick] {0.5*x^2 + 1};
                \addplot[  
                        only marks,  
                        mark=*,  
                        mark size=2pt,  
                        accent  
                    ] coordinates {  
                        (-2,3)  
                        (-1,1.5)  
                        (0,1)  
                        (1,1.5)  
                        (2,3)  
                    };
                \end{axis}
            \end{tikzpicture}
        \end{minipage}
    \end{center}
\end{bsp}

\pagebreak
\subsubsection{Eigenschaften von Funktionen}
\subsubsection*{Stetigkeit}
Eine Funktion $f$ ist stetig wenn sie sprungfrei ist. Das heißt: "Man kann den Graphen der Funktion zeichnen, ohne den Stift abzusetzen".
\begin{center}
  \begin{minipage}{0.45\textwidth}
    \centering
    \begin{tikzpicture}
      \begin{axis}[
        width=7cm, height=5cm,
        axis lines=middle,
        xmin=-2.2, xmax=2.2,
        ymin=-1.1, ymax=3.1,
        samples=200,
        xlabel={$x$}, ylabel={$y$},
        xtick={-2,-1,0,1,2},
        ytick={-1,0,1,2,3},
        domain=-2:2
      ]
        \addplot[accent, ultra thick] {x^3 - x + 1};
      \end{axis}
    \end{tikzpicture}\\[0.5em]
    \captionof*{figure}{stetige Polynomfunktion}
  \end{minipage}\hfill
  \begin{minipage}{0.45\textwidth}
    \centering
    \begin{tikzpicture}
      \begin{axis}[
        width=7cm, height=5cm,
        axis lines=middle,
        xmin=-2.2, xmax=2.2,
        ymin=-1.1, ymax=3.1,
        xlabel={$x$}, ylabel={$y$},
        xtick={-2,-1,0,1,2},
        ytick={-1,0,1,2,3},
        domain=-2:2
      ]
        % unstetige Funktion mit zwei Sprungstellen, z.B. stückweise
        \addplot[accent, ultra thick, domain=-2:-0.5] {1};
        \addplot[accent, ultra thick, domain=-0.5:0.5] {2};
        \addplot[accent, ultra thick, domain=0.5:2] {0};

        % offene/geschlossene Punkte an den Sprungstellen
        % bei x = -0.5
        \addplot[only marks, mark=o, mark size=2pt, accent] coordinates {(-0.5,1)};
        \addplot[only marks, mark=*, mark size=2pt, accent] coordinates {(-0.5,2)};
        % bei x = 0.5
        \addplot[only marks, mark=o, mark size=2pt, accent] coordinates {(0.5,2)};
        \addplot[only marks, mark=*, mark size=2pt, accent] coordinates {(0.5,0)};
      \end{axis}
    \end{tikzpicture}\\[0.5em]
    \captionof*{figure}{Unstetige Funktion mit Sprungstellen}
  \end{minipage}
\end{center}

\vspace{15pt}
\noindent

\begin{minipage}{0.45\textwidth}
{
    \textbf{Monoton steigend (ms)}\\[0.3em]
    $f$ heißt \emph{monoton steigend} auf einem Intervall $I$ , wenn für alle
    $x_1,x_2 \in I$ gilt:
    \begin{equation*}
        \text{wenn } x_1 < x_2 \text{ gilt: } f(x_1) \le f(x_2).
    \end{equation*}
}
\vspace{0.3em}
{
    \centering
    \begin{tikzpicture}
      \begin{axis}[
        width=7cm, height=5cm,
        axis lines=middle,
        xmin=-2.2, xmax=2.2,
        ymin=-1.1, ymax=3.1,
        samples=200,
        xlabel={$x$}, ylabel={$y$},
        xtick={-2,-1,0,1,2},
        ytick={-1,0,1,2,3},
        domain=-2:2
      ]
        % Beispiel: f(x) = x^3 + 1 (streng steigend, aber hier als Beispiel für mon. steigend)
        \addplot[accent, ultra thick] {x^3 + 1};
      \end{axis}
    \end{tikzpicture}
}
\end{minipage}\hfill
\begin{minipage}{0.45\textwidth}
{
    \textbf{Streng monoton steigend (sms)}\\[0.3em]
    $f$ heißt \emph{monoton steigend} auf einem Intervall $I$ , wenn für alle
    $x_1,x_2 \in I$ gilt:
    \begin{equation*}
        \text{wenn } x_1 < x_2 \text{ gilt: } f(x_1) < f(x_2).
    \end{equation*}
}
\vspace{0.3em}
{
    \centering
    \begin{tikzpicture}
      \begin{axis}[
        width=7cm, height=5cm,
        axis lines=middle,
        xmin=-2.2, xmax=2.2,
        ymin=-1.1, ymax=3.1,
        samples=200,
        xlabel={$x$}, ylabel={$y$},
        xtick={-2,-1,0,1,2},
        ytick={-1,0,1,2,3},
        domain=-2:2
      ]
        % Beispiel: f(x) = e^x
        \addplot[accent, ultra thick] {exp(x)/2};
      \end{axis}
    \end{tikzpicture}
}
\end{minipage}

\vspace{1em}

\begin{minipage}{0.45\textwidth}
{
    \textbf{Monoton fallend (mf)}\\[0.3em]
    $f$ heißt \emph{monoton fallend} auf einem Intervall $I$ , wenn für alle
    $x_1,x_2 \in I$ gilt:
    \begin{equation*}
        \text{wenn } x_1 < x_2 \text{ gilt: } f(x_1) \geq f(x_2).
    \end{equation*}
}
\vspace{0.3em}
{
    \centering
    \begin{tikzpicture}
      \begin{axis}[
        width=7cm, height=5cm,
        axis lines=middle,
        xmin=-2.2, xmax=2.2,
        ymin=-1.1, ymax=3.1,
        samples=200,
        xlabel={$x$}, ylabel={$y$},
        xtick={-2,-1,0,1,2},
        ytick={-1,0,1,2,3},
        domain=-2:2
      ]
        % Beispiel: f(x) = -x^2 + 2 (nicht streng, da Scheitelpunkt)
        \addplot[accent, ultra thick] {-x^3 + 1};
      \end{axis}
    \end{tikzpicture}
}
\end{minipage}\hfill
\begin{minipage}{0.45\textwidth}
{
    \textbf{Streng monoton fallend (smf)}\\[0.3em]
    $f$ heißt \emph{streng monoton fallend} auf einem Intervall $I$ , wenn für alle
    $x_1,x_2 \in I$ gilt:
    \begin{equation*}
        \text{wenn } x_1 < x_2 \text{ gilt: } f(x_1) > f(x_2).
    \end{equation*}
}
\vspace{0.3em}
{
    \centering
    \begin{tikzpicture}
      \begin{axis}[
        width=7cm, height=5cm,
        axis lines=middle,
        xmin=-2.2, xmax=2.2,
        ymin=-1.1, ymax=3.1,
        samples=200,
        xlabel={$x$}, ylabel={$y$},
        xtick={-2,-1,0,1,2},
        ytick={-1,0,1,2,3},
        domain=-2:2
      ]
        % Beispiel: f(x) = -e^x
        \addplot[accent, ultra thick] {-exp(x)/2 + 2};
      \end{axis}
    \end{tikzpicture}
}
\end{minipage}

\pagebreak

\subsubsection*{Lokale Extrempunkte, Krümmung und Wendepunkte}
Ein Punkt $x_0$ hei{\ss}t \textbf{lokales Maximum}, wenn $f(x_0)$ \emph{in einer Umgebung} von $x_0$ gr\"o{\ss}er oder gleich allen anderen Funktionswerten ist. \\
Entsprechend hei{\ss}t $x_0$ \textbf{lokales Minimum}, wenn $f(x_0)$ in einer Umgebung von $x_0$ kleiner oder gleich allen anderen Funktionswerten ist. \vspace{5pt} \\
Man sagt, die Funktion ist in einem Intervall \textbf{linksgekrümmt} (positiv gekrümmt, konvex), wenn ihre Tangenten unterhalb des Graphen verlaufen, und \textbf{rechtsgekrümmt} (negativ gekrümmt, konkav), wenn ihre Tangenten oberhalb des Graphen verlaufen.\vspace{5pt} \\
Ein Punkt $x_0$ heißt \textbf{Wendepunkt} einer Funktion $f$, wenn sich dort das Krümmungsverhalten der Funktion ändert, also z.B.\ von „linksgekrümmt“ zu „rechtsgekrümmt“ oder umgekehrt.

\begin{center}  
  \begin{tikzpicture}  
    \begin{axis}[  
      width=0.7\textwidth,  
      axis lines=middle,  
      xmin=-2.2, xmax=2.2,  
      ymin=-1.2, ymax=4.2,  
      samples=200,  
      xlabel={$x$}, ylabel={$y$},  
      xtick={-2,-1,0,1,2},  
      ytick={-1,0,1,2,3,4},  
      domain=-2:2  
    ]  
      % Polynom 3. Grades mit Max und Min, z.B. f(x) = x^3 - 3x  
      \addplot[accent, ultra thick] {x^3 - 3*x + 1};  
  
      % Lokales Maximum bei x = -1  
      \addplot[only marks, mark=*, mark size=2pt, black] coordinates {(-1,3)};  
      \node[above left, black] at (axis cs:-1,3) {Max};  
  
      % Lokales Minimum bei x = 1  
      \addplot[only marks, mark=*, mark size=2pt, black] coordinates {(1,-1)};  
      \node[above=2pt, black] at (axis cs:1,-1) {Min};  

      % Wendepunkt bei x = 0
      \addplot[only marks, mark=*, mark size=2pt, black] coordinates {(0,1)};  
      \node[above right, black] at (axis cs:0,1) {WP};  

      % Krümmungspfeile
      \draw [->, ultra thick] (-1.5,1) .. controls (-1.2,2) and (-0.8,2) .. (-0.5,1);
      \draw [->, ultra thick] (0.5,1) .. controls (0.8,0) and (1.2,0) .. (1.5,1);

      \node[] at (axis cs:-1,1.2) {\textbf{$-$}};
      \node[] at (axis cs:1,0.8) {\textbf{$+$}};
      
    \end{axis}  
  \end{tikzpicture}  
\end{center}

\vspace{10pt}

\begin{center}  
  \begin{minipage}[t]{0.45\textwidth}  
    \textbf{Achsensymmetrie zur $y$-Achse}\\[0.3em]  
    Eine Funktion heißt \emph{achsensymmetrisch zur $y$-Achse}, wenn ihr Graph an der $y$-Achse spiegeln lässt.
    \[  
      f(-x) = f(x)
    \]  
  
    \vspace{0.5em}  
    {
        \centering
        \begin{tikzpicture}  
          \begin{axis}[  
            width=\linewidth,  
            axis lines=middle,  
            xmin=-4.4, xmax=4.4,  
            ymin=-1.4, ymax=5.4,  
            samples=200,  
            xlabel={$x$}, ylabel={$y$},  
            xtick={-4,-2,0,2,4},  
            ytick={-1,0,1,2,3,4,5},  
            domain=-4:4
          ]  
            \addplot[accent, ultra thick] {x^2};    
            \addplot[accent2, ultra thick] {cos(deg(x)) + 2};  
          \end{axis}  
    \end{tikzpicture}  
    }
  \end{minipage}\hfill  
  \begin{minipage}[t]{0.45\textwidth}  
    \textbf{Punktsymmetrie zum Ursprung}\\[0.3em]  
    Eine Funktion heißt \emph{punktsymmetrisch zum Ursprung}, wenn sich ihr Graph durch den Ursprung spiegeln lässt.
    \[  
      f(-x) = -f(x)
    \]  
  
    \vspace{0.5em}  
    {
        \centering
        \begin{tikzpicture}  
          \begin{axis}[  
            width=\linewidth,  
            axis lines=middle,  
            xmin=-4.4, xmax=4.4,  
            ymin=-4.4, ymax=4.4,  
            samples=200,  
            xlabel={$x$}, ylabel={$y$},  
            xtick={-4,-2,0,2,4},  
            ytick={-4,-2,0,2,4},  
            domain=-4:4,  
            legend style={at={(0.02,0.98)},anchor=north west, font=\scriptsize}  
          ]  
            \addplot[accent, ultra thick] {x^3/4};    
            \addplot[accent2, ultra thick] {2*sin(deg(x))};  
          \end{axis}  
        \end{tikzpicture}  
    }
  \end{minipage}  
\end{center}

\subsection{Lineare Funktionen}
Lineare Funktionen sind Funktionen, dereren Graph eine Gerade ist:
\begin{align}
    f(x) = k x + d, \qquad   k &= \text{Steigung} \\
                            d &= \text{Ordinatenabschnitt} \nonumber
\end{align}
Die Steigung $k$ kann über das Steigungsdreieck ermittelt werden:
\begin{equation}
    k = \frac{\Delta y}{\Delta x}
\end{equation}
Mit den trigonometrischen Formeln kann auch der \emph{Steigungswinkel} berechnet werden:
\begin{gather}
    \tan{\alpha} = \frac{\Delta y}{\Delta x} = k \\[5pt]
    \alpha = \arctan{\left( \frac{\Delta y}{\Delta x} \right)} = \arctan{k}
\end{gather}
Der Ordinatenabschnitt $d$ ist der Schnittpunkt der Geraden mit der $y$-Achse.

\begin{bsp}
    Es ist die Lineare Funktion zwischen den zwei Punkten \(A(1,2),\;B(3,3)\) gesucht. \\
    Zuerst wird die Steigung berechnet:
    \begin{align*}
        \Delta y &= B_y - A_y = 3 - 2 = 1 \\
        \Delta x &= B_x - A_x = 3 - 1 = 2 \\[5pt]
        k &= \frac{\Delta y}{\Delta x} = \frac{1}{2} = 0.5
    \end{align*}
    Nachdem das $k$ berechnet wurde kann ein Punkt ($A$ oder $B$) eingesetzt werden um den Ordinatenabschnitt $d$ zu berechnen.
    \begin{align*}
        A_y &= 0.5 \cdot A_x + d \\
        2 &= 0.5 \cdot 1 + d \\
        d &= 1.5
    \end{align*}
    Die Funktionsgleichung lautet also:
    \begin{equation*}
        f(x) = 0.5 x + 1.5
    \end{equation*}
    Der Graph sieht wie folgt aus:

    \begin{center}
    \begin{tikzpicture}
    \begin{axis}[
        axis lines=middle,
        xmin=-1.5, xmax=5,
        ymin=-0.9, ymax=4.4,
        xlabel={$x$},
        ylabel={$y$},
        xtick={-1,0,1,2,3,4},
        ytick={0,1,2,3,4},
        width=10cm,
        height=6.5cm,
        enlargelimits=false,
        axis line style={latex-latex},
    ]
    
        % Gerade f(x) = 0.5x + 1.5
        \addplot[ultra thick, accent] coordinates {(-1,1) (4,3.5)};
        \node[above right] at (axis cs:4,3.5) {\textcolor{accent}{$f(x)$}};
    
        % Punkte A(1,2) und B(3,3)
        \addplot[only marks, mark=*, mark size=2pt] coordinates {(1,2) (3,3)};
        \node[above=4pt] at (axis cs:1,2) {$A(1,2)$};
        \node[above=4pt] at (axis cs:3,3) {$B(3,3)$};
    
        % Steigungsdreieck zwischen A und B
        \addplot[thick] coordinates {(1,2) (3,2)}; % Δx
        \addplot[thick] coordinates {(3,2) (3,3)}; % Δy
    
        % Beschriftungen Δx und Δy
        \node[anchor=north] at (axis cs:2,2) {$\Delta x = 2$};
        \node[anchor=west] at (axis cs:3,2.5) {$\Delta y = 1$};
    
    \end{axis}
    \end{tikzpicture}
    \end{center}
\end{bsp}

\subsection{Potenzfunktionen}
Potenzfunktionen sind Funktionen der Form:
\begin{equation}
    f(x) = ax^z + b, \quad z \in \mathbb{Z}
\end{equation}
Hierbei bestimmt $z$ den Grad der Funktion. Der Parameter $a$ verursacht eine Streckung, Stauchung oder Inversion, und der Parameter $b$ eine Verschiebung in $y$-Richtung. \\
Abhängig von $n$ kann man vier Fälle unterscheiden.
% 1. Zeile: Parabeln
\vspace{14pt}
\\
\noindent
% 1. Zeile, 1. Spalte: n gerade, positiv (x^2, x^4)
\begin{minipage}[t]{0.48\textwidth}
  \centering
  \textbf{Parabeln: $n$ gerade, positiv}
  \[
    x^2,\, x^4,\,\dots
  \]
  \begin{tikzpicture}
    \begin{axis}[
      width=7cm,
      height=5cm,
      axis lines=middle,
      xmin=-2.2, xmax=2.2,
      ymin=-1, ymax=5.5,
      samples=200,
      domain=-2:2,
      xtick={-2,-1,0,1,2},
      ytick={0,1,2,3,4,5},
      axis line style={latex-latex},
    ]
      \addplot[accent, ultra thick] {x^2};
      \addplot[accent2, ultra thick] {x^4};
    \end{axis}
  \end{tikzpicture}
\end{minipage}
\hfill
% 1. Zeile, 2. Spalte: n ungerade, positiv (x^3, x^5)
\begin{minipage}[t]{0.48\textwidth}
  \centering
  \textbf{Parabeln: $n$ ungerade, positiv}
    \[
        x^3,\, x^5,\,\dots
    \]
  \begin{tikzpicture}
    \begin{axis}[
      width=7cm,
      height=5cm,
      axis lines=middle,
      xmin=-2.2, xmax=2.2,
      ymin=-5.5, ymax=5.5,
      samples=200,
      domain=-2:2,
      xtick={-2,-1,0,1,2},
      ytick={-4,-2,0,2,4},
      axis line style={latex-latex},
    ]
      \addplot[accent, ultra thick] {x^3};
      \addplot[accent2, ultra thick] {x^5};
    \end{axis}
  \end{tikzpicture}
\end{minipage}
\vspace{2cm}
% 2. Zeile: Hyperbeln
% 2. Zeile, 1. Spalte: n gerade, negativ (x^{-2}, x^{-4})
\begin{minipage}[t]{0.48\textwidth}
  \centering
  \textbf{Hyperbeln: $n$ gerade, negativ}
  \[
    x^{-2}=\frac{1}{x^2},\,x^{-4}=\frac{1}{x^4},\,\dots
  \]
  \begin{tikzpicture}
    \begin{axis}[
      width=7cm,
      height=5cm,
      axis lines=middle,
      xmin=-3.4, xmax=3.4,
      ymin=-1, ymax=5.5,
      samples=200,
      xtick={-3,-2,-1,0,1,2,3},
      ytick={0,1,2,3,4,5},
      axis line style={latex-latex},
    ]
      % x^{-2}
      \addplot[accent, ultra thick, domain=-3:-0.3] {x^-2};
      \addplot[accent, ultra thick, domain=0.3:3]   {x^-2};
      % x^{-4}
      \addplot[accent2, ultra thick, domain=-3:-0.3] {x^-4};
      \addplot[accent2, ultra thick, domain=0.3:3]   {x^-4};
    \end{axis}
  \end{tikzpicture}
\end{minipage}
\hfill
% 2. Zeile, 2. Spalte: n ungerade, negativ (x^{-3}, x^{-5})
\begin{minipage}[t]{0.48\textwidth}
  \centering
  \textbf{Hyperbeln: $n$ ungerade, negativ}
  \[
    x^{-3}=\frac{1}{x^3},\,x^{-5}=\frac{1}{x^5},\,\dots
  \]
  \begin{tikzpicture}
    \begin{axis}[
      width=7cm,
      height=5cm,
      axis lines=middle,
      xmin=-3.4, xmax=3.4,
      ymin=-5.2, ymax=5.2,
      samples=200,
      xtick={-3,-2,-1,0,1,2,3},
      ytick={-4,-2,0,2,4},
      axis line style={latex-latex},
    ]
      % x^{-3}
      \addplot[accent, ultra thick, domain=-3:-0.3] {x^-3};
      \addplot[accent, ultra thick, domain=0.3:3]   {x^-3};
      % x^{-5}
      \addplot[accent2, ultra thick, domain=-3:-0.3] {x^-5};
      \addplot[accent2, ultra thick, domain=0.3:3]   {x^-5};
    \end{axis}
  \end{tikzpicture}
\end{minipage}
\noindent
\begin{minipage}[t]{0.48\textwidth}
  \centering
  \textbf{Streckung, Stauchung und Spiegelung} \\[0.5em]
  \begin{tikzpicture}
    \begin{axis}[
      width=9cm,
      height=7cm,
      axis lines=middle,
      xmin=-3.4, xmax=3.4,
      ymin=-10, ymax=10,
      samples=200,
      xtick={-3,-2,-1,0,1,2,3},
      ytick={-8,-4,0,4,8},
      axis line style={latex-latex},
    ]
      \addplot[accent, very thick, domain=-3:3] {2*x^2} node[pos=0.75, left] {$2x^2$};
      \addplot[accent2, very thick, domain=-3:3] {x^2} node[pos=0.9, left] {$x^2$};
      \addplot[accent3, very thick, domain=-3:3] {0.5*x^2} node[pos=0.95, below=5pt] {$0.5x^2$};
      \addplot[accent4, very thick, domain=-3:3] {-x^2} node[pos=0.85, right] {$-x^2$};
    \end{axis}
  \end{tikzpicture}
\end{minipage}
\hfill
\begin{minipage}[t]{0.48\textwidth}
  \centering
  \textbf{Verschiebung in $y$-Richtung} \\[0.5em]
  \begin{tikzpicture}
    \begin{axis}[
      width=9cm,
      height=7cm,
      axis lines=middle,
      xmin=-3.4, xmax=3.4,
      ymin=-6, ymax=6,
      samples=200,
      xtick={-3,-2,-1,0,1,2,3},
      ytick={-4,-2,-1,0,1,2,4},
      axis line style={latex-latex},
    ]
      \addplot[accent, very thick, domain=-3:3] {x^2 + 1} node[pos=0.75, left] {$x^2+1$};
      \addplot[accent2, very thick, domain=-3:3] {x^2} node[pos=0.7, right] {$x^2$};
      \addplot[accent3, very thick, domain=-3:3] {0.5*x^2 - 1} node[pos=0.9, below=15pt] {$0.5x^2 - 1$};
      \addplot[accent4, very thick, domain=-3:3] {-x^2 + 0.5} node[pos=0.68, right] {$-x^2 + 0.5$};
    \end{axis}
  \end{tikzpicture}
\end{minipage}

\pagebreak
\subsection{Polynomfunktionen}
Funktionen der Form:
\begin{equation}
    f(x) = a_0 x^0 + a_1 x^1 + a_2 x^2 + \dots + a_n x^n
\end{equation}
heißen Polynomfunktionen $n$-ten Grades. Die höchst vorkommende Potenz bestimmt dabei immer den Grad der Polynomfunktion.\\

\noindent
% Row: 0. Grades
\noindent
\begin{minipage}[c]{0.48\textwidth}
\textbf{Polynomfunktion 0. Grades}
\begin{equation}
    f(x) = a_0 x^0 = a_0
\end{equation}
Polynomfunktionen 0. Grades sind konstante Funktionen. Sie liefern für jeden $x$-Wert den $y$-Wert $a_0$.
\end{minipage}\hfill
\begin{minipage}[c]{0.48\textwidth}
\begin{center}
    \begin{tikzpicture}
        \begin{axis}[
          width=7cm,
          height=5cm,
          axis lines=middle,
          samples=200,
          axis line style={latex-latex},
          xtick={-4,-2,0,2,4},
          ytick={-2,-1,0,1,2,3},
          xmin=-4.5,
          xmax=4.5,
          ymin=-2.5,
          ymax=3.5,
        ]
        \addplot[accent, ultra thick, domain=-3.5:3.5]{2}
        node[pos=0.8, above] {$f_1(x)=2$};
        \addplot[accent2, ultra thick, domain=-3.5:3.5]{-1}
        node[pos=0.8, below] {$f_2(x)=-1$};
        \end{axis}
    \end{tikzpicture}
\end{center}
\end{minipage}

\vspace{2em}

% Row: 1. Grades
\noindent
\begin{minipage}[c]{0.48\textwidth}
\textbf{Polynomfunktion 1. Grades}
\begin{align}
    f(x)&=a_0x^0 + a_1x^1 \\
        &=k x + d \nonumber
\end{align}
Polynomfunktionen 1. Grades sind lineare Funktionen. Sie haben immer genau eine Nullstelle.
\end{minipage}\hfill
\begin{minipage}[c]{0.48\textwidth}
\begin{center}
    \begin{tikzpicture}
        \begin{axis}[
          width=7cm,
          height=5cm,
          axis lines=middle,
          samples=200,
          axis line style={latex-latex},
          xtick={-4,-2,0,2,4},
          ytick={-2,-1,0,1,2,3},
          xmin=-4.5,
          xmax=4.5,
          ymin=-2.5,
          ymax=3.5,
        ]
          \addplot[accent, ultra thick, domain=-3.5:3.5]{0.5*x + 1}
            node[pos=0.8, below=7pt] {$0.5x+1$};
        \end{axis}
    \end{tikzpicture}
\end{center}
\end{minipage}

\vspace{2em}

% Row: 2. Grades
\noindent
\begin{minipage}[c]{0.48\textwidth}
\textbf{Polynomfunktion 2. Grades}
\begin{align}
    f(x) &= a_0 x^0 + a_1 x^1 + a_2 x^2 \\
         &= a x^2 + b x + c \nonumber
\end{align}
Polynomfunktionen 2. Grades sind \emph{quadratische Funktionen} deren Graph eine Parabel ist. Diese werden in weiterer Folge noch genauer untersucht.
\end{minipage}\hfill
\begin{minipage}[c]{0.48\textwidth}
\begin{center}
    \begin{tikzpicture}
        \begin{axis}[
          width=7cm,
          height=5cm,
          axis lines=middle,
          samples=200,
          axis line style={latex-latex},
          xtick={-4,-2,0,2,4},
          ytick={-2,-1,0,1,2,3},
          xmin=-4.5,
          xmax=4.5,
          ymin=-2.5,
          ymax=3.5,
        ]
        \addplot[accent, ultra thick, domain=-4:4]{x^2-x-1};
        \addplot[accent2, ultra thick, domain=-3.8:-0.2]{-0.5*x^2-2*x+-2.5};
        \end{axis}
    \end{tikzpicture}
\end{center}
\end{minipage}

\vspace{2em}

% Row: 3. Grades
\noindent
\begin{minipage}[c]{0.48\textwidth}
\textbf{Polynomfunktion 3. Grades}
\begin{equation}
    f(x) = a_0 x^0 + a_1 x^1 + a_2 x^2 + a_3 x^3
\end{equation}
Polynomfunktionen 3. Grades haben können maximal 3 Nullstellen haben und haben immer zumindest eine.
\end{minipage}\hfill
\begin{minipage}[c]{0.48\textwidth}
\begin{center}
    \begin{tikzpicture}
        \begin{axis}[
          width=7cm,
          height=5cm,
          axis lines=middle,
          samples=200,
          axis line style={latex-latex},
          xtick={-4,-2,0,2,4},
          ytick={-2,-1,0,1,2,3},
          xmin=-4.5,
          xmax=4.5,
          ymin=-2.5,
          ymax=3.5,
        ]
        \addplot[accent, ultra thick, domain=-4:4]{x^3+x^2-2*x-0.5};
        \addplot[accent2, ultra thick, domain=-4:4]{-0.1*x^3+0.4*x^2-0.5*x+1};
        \end{axis}
    \end{tikzpicture}
\end{center}
\end{minipage}

\vspace{2em}

% Row: 4. Grades
\noindent
\begin{minipage}[c]{0.48\textwidth}
\textbf{Polynomfunktion 4. Grades}
\begin{equation}
    f(x) = a_0 x^0 + a_1 x^1 + a_2 x^2 + a_3 x^3 + a_4 x^4
\end{equation}
Polynomfunktionen 4. Grades haben können maximal 4 Nullstellen haben. Es ist aber auch möglich, dass sie gar keine Nullstelle haben.
\end{minipage}\hfill
\begin{minipage}[c]{0.48\textwidth}
\begin{center}
    \begin{tikzpicture}
        \begin{axis}[
          width=7cm,
          height=5cm,
          axis lines=middle,
          samples=200,
          axis line style={latex-latex},
          xtick={-4,-2,0,2,4},
          ytick={-2,-1,0,1,2,3},
          xmin=-4.5,
          xmax=4.5,
          ymin=-2.5,
          ymax=3.5,
        ]
          \addplot[accent, ultra thick, domain=-4:4]{x^4-3*x^3+x^2+x+1};
          \addplot[accent2, ultra thick, domain=-4:4]{-0.6*x^4-1.3*x^3+0.7*x^2+1.6*x-1};
        \end{axis}
    \end{tikzpicture}
\end{center}
\end{minipage}

\subsubsection*{Zusammenhang zwischen Grad und Eigenschaften}
Eine Polynomfunktion $n$-ten Grades kann:
\begin{itemize}
    \item maximal $n$ reelle Nullstellen haben
    \item maximal $n-1$ Extrempunkte haben
    \item maximal $n-2$ Wendepunkte haben
\end{itemize}
Polynomfunktionen bei denen $n$ ungerade ist, haben immer mindestens eine Nullstelle. \\
Polynomfunktionen bei denen $n$ gerade ist, können auch gar keine reelle Nullstelle haben.

\subsubsection*{Quadratische Funktionen}
Quadratische Funktionen sind Polynomfunktionen 2. Grades. Sie werden meist in der allgemeinen Form angeschrieben:
\begin{equation}
    f(x) = ax^2 + bx +c
\end{equation}
Als weitere Form existiert noch die Scheitelpunktform:
\begin{equation}
    f(x) = a (x-x_s)^2 + y_s
\end{equation}
Aus der Scheitelpunktform kann man direkt die Koordinaten des Scheitelpunktes ablesen $S(x_s,y_s)$.
\begin{bsp}
    Gegeben sei die quadratische Funktion in allgemeiner Form:  
    \[  
        f(x) = x^2 + 4x + 3  
    \]  
    Wir wandeln nun in die Scheitelpunktform um (Quadratische Ergänzung):  
    \begin{align*}  
    f(x) &= x^2 + 4x + 3 \\  
         &= \bigl(x^2 + 4x + 4\bigr) - 1 \\  
         &= (x + 2)^2 - 1  
    \end{align*}  
    Damit liegt die Funktion in Scheitelpunktform vor:  
    \[  
        \boxed{f(x) = (x + 2)^2 - 1}
    \]  
    Der Scheitelpunkt hat also die Koordinaten \(S(-2 \mid -1).\). 
    
    \begin{center}  
    \begin{tikzpicture}  
    \begin{axis}[  
        axis lines = middle,  
        xmin = -4.5, xmax = 1.5,  
        ymin = -2.5, ymax = 4.5,  
        xlabel = {$x$},  
        ylabel = {$y$},  
        xtick = {-4,-3,-2,-1,0,1},  
        ytick = {-2,-1,0,1,2,3,4},  
        axis line style={latex-latex}
    ]  
      
    % Parabel: f(x) = (x + 2)^2 - 1  
    \addplot[domain=-5:2, samples=200, accent, ultra thick] {(x + 2)^2 - 1};  
      
    % Scheitelpunkt S(-2 | -1)  
    \addplot[only marks, mark=*, mark size=2pt, black] coordinates {(-2,-1)};  
    \node[below, black] at (axis cs:-2,-1) {$S(-2\mid -1)$};  
      
    \end{axis}  
    \end{tikzpicture}  
    \end{center}
\end{bsp}

\subsection{Exponentialfunktionen}

Funktionen der Form
\begin{equation}
    f(x) = a b^x
\end{equation}
heißen Exponentialfunktionen. Für die Modellierung von kontinuierlichen physikalischen/biologischen Prozessen wird auch oft folgende Form verwendet:
\begin{equation}
    N(t) = N_0 e^{\lambda t}
\end{equation}
Eine wichtige Eigenschaft ist der Schnittpunkt mit der $y$-Achse, welchen man direkt aus beiden Formen ablesen kann. Setzt man $x=0$ bzw. $t=0$ ein erhält man:
\begin{equation}
\begin{array}{c@{\;}c@{\;}c}
    f(0) = a \cdot b^0 = a & \qquad & P(0,a) \\[5pt]
    N(0) = N_0 \cdot e^{\lambda \cdot 0} = N_0 & \qquad & P(0,N_0)
\end{array}
\end{equation}
Die Parameter $b$ und $\lambda$ bestimmen die Steilheit und ob die Exponentialfunktion steigend/fallend ist.\\[5pt]
\noindent
\begin{minipage}[t]{0.45\textwidth}
    \textbf{Lineares Wachstum} \\
    Eine Größe nimmt in gleichen Zeitschritten immer um denselben \emph{absoluten} Betrag zu.
\end{minipage}
\hfill
\begin{minipage}[t]{0.45\textwidth}
    \textbf{Exponentielles Wachstum} \\
    Eine Größe nimmt in gleichen Zeitschritten immer um denselben \emph{prozentualen / relativen} Faktor zu.
\end{minipage}
\vspace{20pt}\\
\begin{minipage}[t]{0.45\textwidth}
    \begin{center}
        \textbf{Wertetabelle von $f(x) = 2x + 1$}\\[0.3cm]
        \renewcommand{\arraystretch}{1.3}  
        \begin{tabular}{|>{\centering\arraybackslash}p{1cm}|>{\centering\arraybackslash}p{1cm}|} 
            \hline
            $x$ & $f(x)$ \\[2pt]
            \hline
            0 & 1  \\
            1 & 3  \\
            2 & 5  \\
            3 & 7  \\
            4 & 9  \\
            5 & 11 \\
            \hline
        \end{tabular}
        \begin{align}
            f(x+1) &= f(x) + k \\
            f(x+n) &= f(x) + k \cdot n
        \end{align}
        \textbf{Graph von $f(x) = 2x + 1$}\\[0.3cm]
        \begin{tikzpicture}
            \begin{axis}[
                width=8cm,
                height=7cm,
                axis lines=middle,
                xlabel={$x$},
                ylabel={$y$},
                xmin=-0.5, xmax=4.5,
                ymin=-1,  ymax=10.5,
                xtick={1,2,3,4},
                ytick={0,2,4,6,8,10},
                grid=both
            ]
            \addplot[domain=0:5,samples=100,very thick,accent]{2*x + 1}
                node[pos=0.6, above=4pt] {$f(x)$};
        
            \draw[thick] (axis cs:0,1) -- (axis cs:1,1);
            \draw[thick] (axis cs:1,1) -- (axis cs:1,3);
            \node[right] at (axis cs:1,2){\small $+2$};
        
            \draw[thick] (axis cs:1,3) -- (axis cs:2,3);
            \draw[thick] (axis cs:2,3) -- (axis cs:2,5);
            \node[right] at (axis cs:2,4){\small $+2$};
        
            \draw[thick] (axis cs:2,5) -- (axis cs:3,5);
            \draw[thick] (axis cs:3,5) -- (axis cs:3,7);
            \node[right] at (axis cs:3,6){\small $+2$};
        
            \draw[thick] (axis cs:3,7) -- (axis cs:4,7);
            \draw[thick] (axis cs:4,7) -- (axis cs:4,9);
            \node[right] at (axis cs:4,8){\small $+2$};
            
            \end{axis}
        \end{tikzpicture}
    \end{center}
\end{minipage}
    \hfill
\begin{minipage}[t]{0.45\textwidth}
    \begin{center}
        \textbf{Wertetabelle von $g(x) = 2\cdot 2^{x}$}\\[0.3cm]
        \renewcommand{\arraystretch}{1.3}  
        \begin{tabular}{|>{\centering\arraybackslash}p{1cm}|>{\centering\arraybackslash}p{1cm}|} 
            \hline
            $x$ & $g(x)$ \\
            \hline
            0 & 2   \\
            1 & 4   \\
            2 & 8   \\
            3 & 16  \\
            4 & 32  \\
            5 & 64  \\
            \hline
        \end{tabular}
        \begin{align}
            f(x+1) &= f(x) \cdot b \\
            f(x+n) &= f(x) \cdot b^n
        \end{align}
        \\
        \textbf{Graph von $g(x) = 2\cdot 2^{x}$}\\[0.3cm]
        \begin{tikzpicture}
            \begin{axis}[
                width=8cm,
                height=7cm,
                axis lines=middle,
                xlabel={$x$},
                ylabel={$y$},
                xmin=-0.5, xmax=4.5,
                ymin=-5,   ymax=55,
                xtick={1,2,3,4},
                ytick={10,20,30,40,50},
                grid=both,
            ]
            \addplot[domain=0:5,samples=100,very thick,accent]{2*2^x}
                node[pos=0.6, above=4pt] {$g(x)$};
            
            \draw[thick] (axis cs:0,2) -- (axis cs:1,2);
            \draw[thick] (axis cs:1,2) -- (axis cs:1,4);
            \node[right] at (axis cs:1,2.45){\small $\cdot 2$};
        
            \draw[thick] (axis cs:1,4) -- (axis cs:2,4);
            \draw[thick] (axis cs:2,4) -- (axis cs:2,8);
            \node[right] at (axis cs:2,6){\small $\cdot 2$};
        
            \draw[thick] (axis cs:2,8) -- (axis cs:3,8);
            \draw[thick] (axis cs:3,8) -- (axis cs:3,16);
            \node[right] at (axis cs:3,12){\small $\cdot 2$};
        
            \draw[thick] (axis cs:3,16) -- (axis cs:4,16);
            \draw[thick] (axis cs:4,16) -- (axis cs:4,32);
            \node[right] at (axis cs:4,24){\small $\cdot 2$};
            
            \end{axis}
        \end{tikzpicture}
    \end{center}
\end{minipage}

\bigskip
\noindent
\begin{minipage}{0.48\textwidth}
{
    \textbf{Wachsende Exponentialfunktionen} \\
    Für $b>1$ und positive $\lambda$ ist die Exponentialfunktion steigend. Je größer der Parameter $b$ ist, desto schneller wächst die Exponentialfunktion. Je größer $\lambda$, desto steiler die Kurve.\\
    \begin{align*}
        f(x) = b^x,\quad b>1\\
        f(x) = e^{\lambda x},\quad \lambda > 0
    \end{align*}
}
{
    \centering
    \begin{tikzpicture}
        \begin{axis}[
            width=\textwidth,
            axis lines=middle,
            xlabel={$x$},
            ylabel={$y$},
            xmin=-2.5, xmax=2.5,
            ymin=-0.5, ymax=6.5,
            xtick = {-2,-1,0,1,2},
            ytick = {0,2,4,6},
            legend style={at={(0.02,0.98)},anchor=north west,font=\small},
            samples=200,
            domain=-2:2,
            axis line style={latex-latex},
        ]
            \addplot[accent,very thick]   {exp(2*x)};
            \addlegendentry{$e^{2x}$}
            
            \addplot[accent2,very thick]  {2^x};
            \addlegendentry{$2^{x}$}
            
            \addplot[accent3,very thick] {4^x};
            \addlegendentry{$4^{x}$}
            
            \addplot[accent4,very thick] {exp(0.5*x)};
            \addlegendentry{$e^{0{,}5x}$}
        \end{axis}
    \end{tikzpicture}
}
\end{minipage}\hfill
\begin{minipage}{0.48\textwidth}
{
    \textbf{Fallende Exponentialfunktionen}
    Ist der Parameter $b$ zwischen 0 und 1 oder der Parameter $\lambda$ negativ, so ist die Exponentialfunktion fallend. Beachte: Ein negativer Exponent mit $b>1$ ergibt auch eine fallend Exponentialfunktion: \(4^{-x}=\frac{1}{4^x}\)
    \begin{align*}
        f(x) = b^x,\quad 0<b<1\\
        f(x) = e^{\lambda x},\quad \lambda < 0
    \end{align*}

}
{
    \centering
    \begin{tikzpicture}
        \begin{axis}[
            width=\textwidth,
            axis lines=middle,
            xlabel={$x$},
            ylabel={$y$},
            xmin=-2.5, xmax=2.5,
            ymin=-0.5, ymax=6.5,
            xtick = {-2,-1,0,1,2},
            ytick = {0,2,4,6},
            legend style={at={(0.98,0.98)},anchor=north east,font=\small},
            samples=200,
            domain=-2:2,
            axis line style={latex-latex},
        ]
            \addplot[accent, very thick]  {2^(-x)};
            \addlegendentry{$2^{-x}$}
            
            \addplot[accent2, very thick] {4^(-x)};
            \addlegendentry{$4^{-x}$}
            
            \addplot[accent3, very thick]   {exp(-2*x)};
            \addlegendentry{$e^{-2x}$}
            
            \addplot[accent4, very thick] {exp(-0.5*x)};
            \addlegendentry{$e^{-0{,}5x}$}
        \end{axis}
    \end{tikzpicture}
}
\end{minipage}

\vspace{20pt}

\noindent
\begin{minipage}[t]{0.48\textwidth}
{
    \textbf{Streckung, Stauchung, Inversion}\\
    Der Faktor $a,\,N_0$ bewirkt eine Streckung \\($a>1$), Stauchung ($0<a<1$) oder Inversion ($a<0$). Er gibt auch den Schnittpunkt mit der $y$-Achse $P(0,a)$ $P(0,N_0)$ an.
}
{
    \begin{align*}
        f(x) &= a \, b^x, \quad b>1 \\
        f(x) &= N_0 \, e^x
    \end{align*}
    \begin{tikzpicture}
      \begin{axis}[
        width=\linewidth,
        height=7cm,
        axis lines=middle,
        xmin=-2.5, xmax=3.5,
        ymin=-4.5, ymax=4.5,
        xtick={-2,-1,0,1,2,3},
        ytick={-4,-3,-2,-1,0,1,2,3,4},
        samples=200,
        domain=-2:3,
        legend style={font=\scriptsize,at={(0.97,0.03)},anchor=south east},
        axis line style={latex-latex},
        xlabel={$x$}, ylabel={$y$}
      ]
        \addplot[accent, very thick]   {2*2^x};
        \addlegendentry{$2\cdot 2^x$}
    
        \addplot[accent2, very thick]    {2^x};
        \addlegendentry{$2^x$}
    
        \addplot[accent3, very thick] {0.5*2^x};
        \addlegendentry{$0{,}5\cdot 2^x$}
    
        \addplot[accent4, very thick] {-2^x};
        \addlegendentry{$-2^x$}
      \end{axis}
    \end{tikzpicture}
}
\end{minipage}\hfill
\begin{minipage}[t]{0.48\textwidth}
{
    \textbf{Verschiebung in $y$‑Richtung}\\
    Die additive Konstanten c bewirkt wie bei allen bisherigen Funktionstypen eine Verschiebung in $y$-Richtung.\\
    \begin{align*}
      f(x) = a\, b^x + c \\
      f(x) = N_0 e^{\lambda x} + c
    \end{align*}
}
{
    \centering
    \begin{tikzpicture}
      \begin{axis}[
        width=\linewidth,
        height=7cm,
        axis lines=middle,
        xmin=-2.5, xmax=3.5,
        ymin=-4.5, ymax=4.5,
        xtick={-2,-1,0,1,2,3},
        ytick={-4,-3,-2,-1,0,1,2,3,4},
        samples=200,
        domain=-2:3,
        legend style={font=\scriptsize,at={(0.97,0.03)},anchor=south east},
        axis line style={latex-latex},
        xlabel={$x$}, ylabel={$y$},
      ]
        \addplot[accent, very thick]   {2^x + 1};
        \addlegendentry{$2^x+1$}
    
        \addplot[accent2, very thick]    {2^x};
        \addlegendentry{$2^x$}
    
        \addplot[accent3, very thick] {2^x - 2};
        \addlegendentry{$2^x-2$}
    
        \addplot[accent4, very thick] {-2^x + 1};
        \addlegendentry{$-2^x+1$}
      \end{axis}
    \end{tikzpicture}
}
\end{minipage}

\pagebreak

\subsubsection{natürliche Wachstumsvorgänge}
Natürliche Wachstumsvorgänge werden meist mit folgender Formel beschrieben:
\begin{align}
    N(t)=N_0 \cdot e^{\lambda t},\quad \lambda > 0 \qquad  &N(t) \dots \text{Größe zum Zeitpunkt } t \\
                                                           &N_0 \dots \text{Anfangsgröße} \nonumber \\
                                                           &\lambda \dots \text{Wachstumskonstante} \nonumber
\end{align}
Ein wichtiger Parameter ist die \emph{Verdopplungszeit}. Sie wird als die Zeit definiert, nach der sich die Grundmenge $N_0$ verdoppelt hat.\\
Um die Verdopplungszeit auszurechnen muss man folgende Gleichung nach $t$ Lösen:
\begin{equation}
    2 N_0 = N_0 \cdot e^{\lambda t}
\end{equation}

\subsubsection{natürliche Abklingvorgänge}
Natürliche Abklingvorgänge werden meist mit folgender Formel beschrieben:
\begin{align}
    N(t)=N_0 \cdot e^{-\lambda t},\quad \lambda > 0 \qquad  &N(t) \dots \text{Größe zum Zeitpunkt } t \\
                                                           &N_0 \dots \text{Anfangsgröße} \nonumber \\
                                                           &\lambda \dots \text{Wachstumskonstante} \nonumber
\end{align}
Ein wichtiger Parameter ist die \emph{Halbwertszeit}. Sie wird als die Zeit definiert, nach der sich die Grundmenge $N_0$ halbiert hat.\\
Um die Halbwertszeit auszurechnen muss man folgende Gleichung nach $t$ Lösen:
\begin{equation}
    \frac{N_0}{2} = N_0 \cdot e^{\lambda t}
    \label{eq:halbwertszeit}
\end{equation}
\vspace{20pt}
\begin{aufgabe}
    \emph{Radioaktiver Zerfall von Arsen-73} \\[2pt]
    Die Anzahl der Atomkerne wird durch $N(t)=N_0 \cdot e^{-\lambda t}$ beschrieben wobei\\
    $N_0=10^{20},\quad \lambda = 0.0086\,d^{-1}$
    \begin{itemize}
        \item[1.] Ermittle die Anzahl an Atomkernen nach 200 Tagen
        \item[2.] Ermittle die Halbwertszeit
        \item[3.] Stelle $N(t)$ graphisch dar und markiere die Halbwertszeit
    \end{itemize}
\end{aufgabe}
\begin{loesung}
    \begin{itemize}
        \item[1.] Die Anzahl erhält man wenn man $200\,d$ als Argument für $N(t)$ verwendet:
            \[
                N(200\,d) = N_0 \cdot e^{-\lambda t} = 10^{20} \cdot e^{-0.0086 \cdot 200} \approx 
                \boxed{1.79 \cdot 10^{19}}
            \]
        \item[2.] Die Halbwertszeit erhalten wir indem wir Gleichung \ref{eq:halbwertszeit} lösen:
            \begin{align*}
                \frac{N_0}{2} &= N_0 \cdot e^{-\lambda t} \\
                \frac{1}{2} &= e^{-\lambda t} \\
                \ln{\left(\frac{1}{2}\right)} &= -\lambda t \ln{(e)} \\
                t &= \frac{\ln{\left(\frac{1}{2}\right)}}{\lambda} \approx \boxed{80.6\,d}
            \end{align*}
            \begin{geogebracode}
                NLöse(1/2=exp(-0.0086*t),t) // CAS            
            \end{geogebracode}
        \item[3.] Die Funktion sieht folgendermaßen aus
    \end{itemize}

    \begin{center}
    \begin{tikzpicture}
        \begin{axis}[
            width=12cm,
            height=8cm,
            axis lines=middle,  
            xlabel={$t\,[\mathrm{d}]$},
            ylabel={$N(t)$},
            xmin=0, xmax=250,
            ymin=0, ymax=1.05e20,
            ytick={2e19,4e19,6e19,8e19,1e20},
            yticklabels={$2\cdot10^{19}$,$4\cdot10^{19}$,$6\cdot10^{19}$,$8\cdot10^{19}$,$10^{20}$},
            scaled y ticks=false,
        ]
            % Funktion N(t) = 10^20 * e^{-0.0086 t} (einzige farbige Kurve)
            \addplot[domain=0:250, samples=200, accent, thick] {1e20 * exp(-0.0086*x)};
            
            % Punkt 1: t = 80.6 d, N = 10^20 / 2 = 5e19
            \addplot[only marks, mark=*, mark size=2pt, black] coordinates {(80.6,5e19)};
            \node[anchor=south west] at (axis cs:80.6,5e19)
              {$\bigl(80.6\,\mathrm{d},\tfrac{10^{20}}{2}\bigr)$};
            \addplot[black, densely dashed] coordinates {(80.6,0) (80.6,5e19)};
            \addplot[black, densely dashed] coordinates {(0,5e19) (80.6,5e19)};
            
            % Punkt 2: t = 200 d, N = 1.77 * 10^19
            \addplot[only marks, mark=*, mark size=2pt, black] coordinates {(200,1.77e19)};
            \node[above] at (axis cs:200,1.77e19)
              {$\bigl(200\,\mathrm{d},1.77\cdot10^{19}\bigr)$};
            \addplot[black, densely dashed] coordinates {(200,0) (200,1.77e19)};
            \addplot[black, densely dashed] coordinates {(0,1.77e19) (200,1.77e19)};
        \end{axis}
    \end{tikzpicture}
    \end{center}
\end{loesung}

\pagebreak
\subsection{Sinusfunktion, Cosinusfunktion}
\subsubsection{Graph der Sinusfunktion}
Lässt man einen Punkt entlang des Einheitskreises laufen, dann entspricht die $y$-Position des Punktes den Wert der Sinusfunktion.
\begin{center}
    \begin{tikzpicture}[scale=1.5]
        \begin{scope} % Einheitskreis
            \draw[thick] (0.4,0) arc (0:45:0.4);    % Winkel
            \node at (0.5,0.2) {$\alpha$};
            
            \coordinate (P) at (45:1);              % Punkt am Einheitskreis
            \draw[thick] (0,0) -- (P);              % Radius
            \draw[thick, accent2] (0.707,0) -- (P); % y-Länge: sin(x)
            \fill (P) circle (1pt);                 % Punkt am Einheitskreis
        
            \draw[thick] (0,0) circle (1);          % Kreis
            \draw[<->, thick] (-1.4,0) -- (1.4,0);  % x-Achse
            \draw[<->, thick] (0,-1.4) -- (0,1.4);  % y-Achse
            
            \node[below right] at (1,0) {$1$};
            \node[above right] at (0,1) {$1$};
            \node[below left] at (-1,0) {$-1$};
            \node[below right] at (0,-1) {$-1$};
        \end{scope}
        
        \coordinate (Psin) at ($(2cm,0) + (0.785,{sin(45)})$);  
        \draw[accent2, thick] ($(2cm,0) + (0.785,0)$) -- (Psin);
        \draw[densely dashed] (P) -- (Psin);
        
        \begin{scope}[xshift=2cm] % Funktionsgraph
            \draw[<->, thick] (-0.2,0) -- (7,0) node[right] {$\alpha$};
            \draw[<->, thick] (0,-1.3) -- (0,1.3) node[above] {$y$};
            
            % Sinusfunktion
            \draw[accent, ultra thick, domain=0:6.283,smooth,variable=\x] plot
            ({\x},{sin(\x r)});
            
            % Markierungen y-Achse 
            \draw[thick] (-0.05,1) -- (0.05,1);
            \node[left] at  (-0.05,1) {$1$};
            \draw[thick] (-0.05,-1) -- (0.05,-1);
            \node[left] at (-0.05,-1) {$-1$};
            
            % Winkelbeschriftungen
            \foreach \ang/\x in {45/0.785,90/1.571,180/3.142,270/4.712,360/6.283} {
            \draw[very thick] (\x,0.05) -- (\x,-0.05);
            \node[below] at (\x,0) {$\ang^\circ$};
            }
        \end{scope}
        \fill (Psin) circle (1pt);
    \end{tikzpicture}
\end{center}
\begin{itemize}
    \item ungerade Funktion (punktsymmetrisch zum Ursprung) 
    \item $2\pi$-periodisch: $\sin{(x+2\pi)} = \sin{(x)}$
    \item Definitionsbereich $D=\mathbb{R}$, Wertebereich $W=[-1,1]$
    \item Nullstellen $x=\{ \dots,\,-2\pi,\,-\pi,\,0,\,\pi,\,2\pi,\,\dots \}$ 
\end{itemize}

\subsubsection{Graph der Cosinusfunktion}
Lässt man einen Punkt entlang des Einheitskreises laufen, dann entspricht die $x$-Koordinate des Punktes den Wert der Sinusfunktion.
\begin{center}
    \begin{tikzpicture}[scale=1.5]
        \begin{scope} % Einheitskreis
            \draw[thick] (0.4,0) arc (0:45:0.4);    % Winkel
            \node at (0.5,0.2) {$\alpha$};
            
            \coordinate (P) at (45:1);              % Punkt am Einheitskreis
            \draw[thick] (0,0) -- (P);              % Radius
            \draw[thick] (0.707,0) -- (P);          % y-Länge: sin(x)
            \fill (P) circle (1pt);                 % Punkt am Einheitskreis
        
            \draw[thick] (0,0) circle (1);          % Kreis
            \draw[<->, thick] (-1.4,0) -- (1.4,0);  % x-Achse
            \draw[<->, thick] (0,-1.4) -- (0,1.4);  % y-Achse
            
            \node[below right] at (1,0) {$1$};
            \node[above right] at (0,1) {$1$};
            \node[below left] at (-1,0) {$-1$};
            \node[below right] at (0,-1) {$-1$};
    
            \draw[thick, accent2] (0,0) -- (0.707,0);
            
        \end{scope} 
        
        \coordinate (Pcos) at ($(2cm,0) + (0.785,{sin(45)})$);  
        \draw[accent2, thick] ($(2cm,0) + (0.785,0)$) -- (Pcos);
    
        \begin{scope}[xshift=2cm] % Funktionsgraph
            \draw[<->, thick] (-0.2,0) -- (7,0) node[right] {$\alpha$};
            \draw[<->, thick] (0,-1.3) -- (0,1.3) node[above] {$y$};
            
            % Cosinusfunktion
            \draw[accent, ultra thick, domain=0:6.283,smooth,variable=\x] plot
            ({\x},{cos(\x r)});
            
            % Markierungen y-Achse 
            \draw[thick] (-0.05,1) -- (0.05,1);
            \node[left] at  (-0.05,1) {$1$};
            \draw[thick] (-0.05,-1) -- (0.05,-1);
            \node[left] at (-0.05,-1) {$-1$};
            
            % Winkelbeschriftungen
            \foreach \ang/\x in {45/0.785,90/1.571,180/3.142,270/4.712,360/6.283} {
            \draw[very thick] (\x,0.05) -- (\x,-0.05);
            \node[below] at (\x,0) {$\ang^\circ$};
            }
        \end{scope}
        \fill (Pcos) circle (1pt);
    \end{tikzpicture}
\end{center}
\begin{itemize}
    \item gerade Funktion (symmetrisch zur $y$-Achse) 
    \item $2\pi$-periodisch: $\cos{(x+2\pi)} = \cos{(x)}$
    \item Definitionsbereich $D=\mathbb{R}$, Wertebereich $W=[-1,1]$
    \item Nullstellen $x = \{\ldots, -\tfrac{3\pi}{2},\, -\tfrac{\pi}{2},\, \tfrac{\pi}{2},\, \tfrac{3\pi}{2},\, \tfrac{5\pi}{2},\, \ldots\}$
\end{itemize}

\subsubsection{Graph der Tangensfunktion}
\begin{center}
    \begin{tikzpicture}[scale=1.5]
        \begin{scope}
            \draw[thick] (0.4,0) arc (0:45:0.4);    % Winkel
            \node at (0.5,0.2) {$\alpha$};
            
            \coordinate (P) at (45:1);              % Punkt am Einheitskreis
            \draw[thick] (0,0) -- (P);              % Radius
            \draw[thick] (0.707,0) -- (P);          % y-Länge: sin(x)
        
            \draw[thick] (0,0) circle (1);          % Kreis
            \draw[<->, thick] (-1.4,0) -- (1.4,0);  % x-Achse
            \draw[<->, thick] (0,-1.4) -- (0,1.4);  % y-Achse
            
            \node[below right] at (1,0) {$1$};
            \node[above right] at (0,1) {$1$};
            \node[below left] at (-1,0) {$-1$};
            \node[below right] at (0,-1) {$-1$};
        
            \draw[thick] (0,0) -- (P);
            \draw[densely dashed] (1,-1.5) -- (1,0);
            \draw[densely dashed] (1,1) -- (1,2);
            \coordinate (T) at (1,{tan(45)});
            \draw[dashed] (P) -- (T);  % Hilfslinie von P zu T
            \draw[thick, accent2] (1,0) -- (T);
            \fill (T) circle (1pt);
        \end{scope}
        
        \draw[densely dashed] (1,1) -- ($(2cm,0) + (0.785,1)$);
        \draw[accent2, thick] ($(2cm,0) + (0.785,0)$) -- ($(2cm,0) + (0.785,1)$);
        
        \begin{scope}[xshift=2cm] % Funktionsraph
            \draw[<->, thick] (-0.2,0) -- (7,0) node[right] {$\alpha$};
            \draw[<->, thick] (0,-1.3) -- (0,1.3) node[above] {$y$};
        
            % Vertikale Asymptoten bei 90° (pi/2) und 270° (3pi/2)
            \draw[densely dashed] (1.571,-2) -- (1.571,2);
            \draw[densely dashed] (4.712,-2) -- (4.712,2);
            
            \clip (-0.5,-2) rectangle (7,2);
        
            % Tangens-Äste: 0°..90°, 90°..270°, 270°..360°
            % 1. Ast: 0 bis etwas vor 90°
            \draw[accent, ultra thick, domain=0.0:1.4, smooth, variable=\x]
              plot ({\x},{tan(\x r)});
            % 2. Ast: etwas nach 90° bis etwas vor 270°
            \draw[accent, ultra thick, domain=1.75:4.5, smooth, variable=\x]
              plot ({\x},{tan(\x r)});
            % 3. Ast: etwas nach 270° bis 360°
            \draw[accent, ultra thick, domain=4.9:6.283, smooth, variable=\x]
              plot ({\x},{tan(\x r)});
        
            % Markierungen y-Achse 
            \draw[thick] (-0.05,1) -- (0.05,1);
            \node[left] at  (-0.05,1) {$1$};
            \draw[thick] (-0.05,-1) -- (0.05,-1);
            \node[left] at (-0.05,-1) {$-1$};
        
            % Winkelbeschriftungen
            \foreach \ang/\x in {45/0.785,90/1.571,135/2.356,180/3.142,225/3.927,270/4.712,315/5.498,360/6.283} {
            \draw[very thick] (\x,0.05) -- (\x,-0.05);
            \node[below] at (\x,0) {$\ang^\circ$};
            }
        \end{scope}
        \fill ($(2cm,0) + (0.785,1)$) circle (1pt);
    \end{tikzpicture}
\end{center}
\begin{itemize}
    \item ungerade Funktion (punktsymmetrisch zum Ursprung) 
    \item $\pi$-periodisch: ($\tan{(x+\pi} = \tan{(x)}$
    \item Definitionsbereich $D= \mathbb{R} \setminus \{ \frac{\pi}{2} + k\pi,\; k \in \mathbb{Z} \}$, Wertebereich $W=\mathbb{R}$
    \item Nullstellen $x = \{\ldots, -2\pi,\, -\pi,\, 0,\, \pi,\, 2\pi,\, \ldots\}$
\end{itemize}

\subsubsection{Allgemeine Sinusfunktion}
Wir betrachten Funktionen der Form
\begin{equation}
    f(x) = a \cdot \sin(bx) + c.
\end{equation}

\paragraph{Bedeutung der Parameter}
\begin{itemize}
  \item[$a:$] Amplitude: Streckung/Stauchung in $y$-Richtung und ggf. Spiegelung an der $x$.
    \item[$b:$] Kreisfrequenz: Streckung/Stauchung in $x$-Richtung, hat einen direkten Zusammenhang mit der Frequenz $f$ und der Periodendauer $T$:
        \begin{align}
            b = 2 \pi \cdot f, \qquad f = \frac{b}{2\pi} \\
            T = \frac{1}{f}=\frac{2\pi}{b}, \qquad b = \frac{2\pi}{T}
        \end{align}
  \item[$c:$]  Additive Konstante: Verschiebung in $y$-Richtung.
\end{itemize}
\noindent
Oft sieht man in technischen Anwendungen auch die Form:
\begin{align*}
    x(t) = \hat{x} \cdot \sin{(\omega \cdot t + \varphi)} + \Bar{x}, 
    \qquad  & \hat{a} \dots  \text{Amplitude}          \\
            & \omega  \dots  \text{Kreisfrequenz}      \\
            & \varphi \dots  \text{Phasenverschiebung} \\
            & \Bar{x} \dots  \text{Mittelwert/Gleichanteil}
\end{align*}
Hier entprechen die Parameter $\hat{x},\omega \text{ und } \Bar{x}$ den Parametern $a,b \text{ und } c$.
Der Parameter $\varphi$ entspricht einer Verschiebung in $x$-Richtung. \\[10pt]
\begin{aufgabe}
    Im folgenden Graph ist eine Sinusfunktion der Form
    \[
        f(x) = a \cdot \sin{(bx)} + d
    \]
    dargestellt. Ermittle die Parameter $a,b \text{ und } c$.

    \begin{center}
      \begin{tikzpicture}
        \begin{axis}[
          width=14cm,
          height=7cm,
          axis lines=middle,
          xmin=-4.5*pi, xmax=4.5*pi,
          ymin=-3, ymax=5,
          xtick={-4*pi,-3*pi,-2*pi,-pi,0,pi,2*pi,3*pi,4*pi},
          xticklabels={$-4\pi$,$-3\pi$,$-2\pi$,$-\pi$,$0$,$\pi$,$2\pi$,$3\pi$,$4\pi$},
          ytick={-2,-1,0,1,2,3,4},
          trig format=rad,   
          xlabel={$x$},
          ylabel={$y$},
          axis line style = {latex-latex},
        ]
    
          % eigentliche Sinusfunktion
          \addplot[
            domain=-4*pi:4*pi,
            samples=300,
            accent,
            very thick
          ] {2*sin(0.5*x) + 1}
                node[pos=0.2, above right] {$f(x)$};    
    
        \end{axis}
      \end{tikzpicture}
    \end{center}
\end{aufgabe}


\begin{loesung} \\
    Zuerst ist es immer hilfreich, die Mittellinie einzuzeichnen
    \begin{center}
      \begin{tikzpicture}[>=stealth]
        \begin{axis}[
          width=14cm,
          height=7cm,
          axis lines=middle,
          xmin=-4.5*pi, xmax=4.5*pi,
          ymin=-3, ymax=5,
          xtick={-4*pi,-3*pi,-2*pi,-pi,0,pi,2*pi,3*pi,4*pi},
          xticklabels={$-4\pi$,$-3\pi$,$-2\pi$,$-\pi$,$0$,$\pi$,$2\pi$,$3\pi$,$4\pi$},
          ytick={-2,-1,0,1,2,3,4},
          trig format=rad,   
          xlabel={$x$},
          ylabel={$y$},
          axis line style = {latex-latex},
        ]
    
          % eigentliche Sinusfunktion
          \addplot[
            domain=-4.1*pi:4.1*pi,
            samples=300,
            accent,
            very thick
          ] {2*sin(0.5*x) + 1}
                node[pos=0.2, above right] {$f(x)$};    

            \addplot[dashed, domain=-4.2*pi:4.2*pi] {1};
                
            \addplot[mark=*] coordinates {(pi,3)} node[above right] {Max};
            \addplot[mark=*] coordinates {(-pi,-1)} node[below left] {Min};
            
            \addplot[mark=*] coordinates {(0,1)};
            
            %\draw[<->] (axis cs:pi,1) -- (axis cs:pi,3) {$a$};
            \draw[<->] (axis cs:pi,1) -- node[right] {$a$} (axis cs:pi,3);
            \draw[<->] (axis cs:pi/2,0) -- node[right] {$c$} (axis cs:pi/2,1);
            \draw[<->] (axis cs:-pi, -1.7) -- node[below] {$T$} (axis cs:3*pi, -1.7);
    
        \end{axis}
      \end{tikzpicture}
    \end{center}
    Die Amplitude $a$ kann man entweder direkt ablesen oder über das Maximum und Minimum berechnen:
    \[ 
        a = \frac{\text{Max}-\text{Min}}{2} = \frac{3-(-1)}{2} = 2
    \]
    Die Frequenz berechnet sich über die Periodendauer welche man als Abstand der zwei Minima bzw. Maxima ablesen kann.
    \[
        b = \frac{2\pi}{T} = \frac{2\pi}{4\pi} = \frac{1}{2}
    \]
    Die Additive Konstante kann man als Abstand zwischen $x$-Achse und Mittellinie ablesen.
    \[
        c = 1
    \]
    Wir erhalten als Lösung
    \[
        \boxed{f(x) =  2 \cdot \sin{\left( \frac{1}{2}\cdot x \right)} + 1}
    \]
\end{loesung}